\documentclass[11pt,a4paper]{article}

% --- Encoding & language -----------------------------------------------------
\usepackage[utf8]{inputenc}
\usepackage[T1]{fontenc}
\usepackage[brazilian]{babel}
\usepackage{lmodern}

% --- Layout ------------------------------------------------------------------
\usepackage[margin=2.3cm]{geometry}
\usepackage{microtype}
\usepackage{parskip}

% --- Tabelas e listas --------------------------------------------------------
\usepackage{booktabs}
\usepackage{array}
\usepackage{tabularx}
\usepackage{enumitem}

% --- Cor, código, gráficos ---------------------------------------------------
\usepackage{xcolor}
\usepackage{listings}
\usepackage{graphicx}
\usepackage{tikz}
\usetikzlibrary{positioning,arrows.meta,shapes.geometric,fit,backgrounds}

% --- Hyperref no final --------------------------------------------------------
\usepackage[
    colorlinks=true,
    linkcolor=black,
    urlcolor=blue!60!black,
    citecolor=black
]{hyperref}

% --- Estilo de código -------------------------------------------------------
\definecolor{codebg}{rgb}{0.96,0.96,0.97}
\definecolor{codekw}{rgb}{0.16,0.30,0.65}
\definecolor{codestr}{rgb}{0.10,0.45,0.20}
\definecolor{codecomment}{rgb}{0.50,0.50,0.55}

\lstset{
    basicstyle=\ttfamily\small,
    backgroundcolor=\color{codebg},
    keywordstyle=\color{codekw}\bfseries,
    stringstyle=\color{codestr},
    commentstyle=\color{codecomment}\itshape,
    showstringspaces=false,
    breaklines=true,
    frame=single,
    rulecolor=\color{black!10},
    framesep=4pt,
    xleftmargin=4pt,
    captionpos=b,
}

% --- Atalhos -----------------------------------------------------------------
\newcommand{\code}[1]{\texttt{\small #1}}
\newcommand{\file}[1]{\texttt{\small\color{black!75} #1}}

% --- Meta --------------------------------------------------------------------
\title{\textbf{Protocolo de Benchmark Integrado Retrospectivo}\\
  \large Pipeline SINAM + ChatTime --- avaliação ponta a ponta de LLMs
  orquestradores e modelos de série temporal em prompts de análise do passado}
\author{Projeto SINAM/VIOLBR \\ \small \texttt{hernandison@gmail.com}}
\date{Junho de 2026}

% =============================================================================
\begin{document}
\maketitle

\begin{abstract}
Este documento descreve o protocolo de avaliação aplicado ao pipeline
SINAM + ChatTime --- uma arquitetura que combina um \emph{LLM orquestrador}
(Qwen3 via Ollama) com modelos de \emph{série temporal} (Chronos--2,
ChatTime, ETS, Random Forest, entre outros) para responder perguntas em
português sobre notificações de violência (sistema VIOLBR). O caso de uso
principal do produto é a \textbf{análise retrospectiva}: o usuário
pergunta \emph{como evoluiu} um tipo de violência em uma UF/município
nos últimos anos, e o sistema responde com tendências, sazonalidade e
diagnóstico estrutural. O benchmark é estruturado em três modalidades:
\textbf{Orquestrador} (avalia só o LLM, com forecaster fixo),
\textbf{Retrospectivo puro} (\emph{walk-forward backtest} dos forecasters
em séries isoladas) e \textbf{Integrado Retrospectivo} (avalia a cadeia
ponta a ponta, com prompts típicos de análise do passado). Esta última
--- o foco deste documento --- faz produto cartesiano
\(\text{LLM} \times \text{forecaster}\) sobre o gold-set de perguntas
retrospectivas e mede sete critérios booleanos por execução, agregados
num \emph{score} ponderado.
\end{abstract}

\tableofcontents
\vspace{1em}

% =============================================================================
\section{Motivação}

\subsection{O caso de uso real é retrospectivo}

O produto SINAM + ChatTime não foi pensado para projetar o futuro --- é
uma ferramenta de \textbf{análise do passado}. Perguntas típicas:

\begin{itemize}[leftmargin=*]
\item \emph{``Como evoluiu a violência sexual em Sergipe nos últimos anos?''}
\item \emph{``Qual a tendência de tortura no Brasil entre 2020 e 2024?''}
\item \emph{``A negligência contra idosos no RJ está crescendo?''}
\end{itemize}

Todas pedem \emph{olhar para trás}: tendência, sazonalidade, mudança de
regime, diagnóstico estrutural. Mesmo quando o pipeline interno usa
forecasters para \emph{cross-validar} a série (backtest em janela móvel),
o que o usuário recebe é uma narrativa retrospectiva.

\subsection{Por que avaliar a cadeia inteira}

A avaliação tradicional separa LLM de modelo preditivo: mede \emph{tool
calling} com um \emph{dataset} sintético e mede acurácia do \emph{forecaster}
em séries isoladas. Em produção, porém, o resultado depende dos dois
funcionarem \emph{juntos}. Se o LLM roteia errado (chama
\code{relatorio\_municipio} pra ``Aracaju'' sem a UF), nem o melhor
forecaster do mundo salva a resposta.

A modalidade \textbf{Integrado Retrospectivo} fecha essa lacuna: ela
responde com dados a pergunta de produto \emph{``qual combinação
\(\langle \text{LLM}, \text{forecaster} \rangle\) vale a pena pro caso
de uso retrospectivo nesse domínio?''}, em vez de tratar cada componente
como otimizável isoladamente.

% =============================================================================
\section{Arquitetura de avaliação}

O sistema é dividido em três \emph{posições} funcionalmente distintas
(figura \ref{fig:arch}). Um mesmo modelo \emph{pode} ocupar mais de uma
posição, mas a avaliação trata cada \emph{slot} separadamente.

\begin{figure}[h!]
\centering
\begin{tikzpicture}[
    node distance=0.6cm and 0.7cm,
    box/.style={
        rectangle, draw=black!40, rounded corners=2pt,
        minimum width=3.6cm, minimum height=1.7cm, align=center,
        font=\small
    },
    arrow/.style={-{Stealth[length=2mm]}, thick, black!60},
    label/.style={font=\footnotesize\itshape, black!60}
]
\node[box, fill=blue!8]  (orq)    {\textbf{Orquestrador}\\(P1)\\Qwen3--8B + tools};
\node[box, fill=green!8, right=of orq] (retro)
    {\textbf{Retrospectivo}\\(P2)\\Explica o passado\\ChatTime, RF, ETS};
\node[box, fill=orange!10, right=of retro] (prosp)
    {\textbf{Prospectivo}\\(P3)\\Projeta o futuro\\Chronos--2, TimesFM};
\draw[arrow] (orq) -- (retro);
\draw[arrow] (retro) -- (prosp);
\node[label, below=0.05cm of orq]   {pergunta PT-BR};
\node[label, below=0.05cm of retro] {séries SINAM};
\node[label, below=0.05cm of prosp] {horizonte futuro};
\end{tikzpicture}
\caption{As três posições da arquitetura. O Integrado mede a cadeia
P1 $\to$ P2/P3 funcionando junta.}
\label{fig:arch}
\end{figure}

% =============================================================================
\section{As três modalidades de \emph{benchmark}}

\begin{table}[h!]
\centering
\small
\begin{tabularx}{\textwidth}{@{}l X l@{}}
\toprule
\textbf{Modalidade} & \textbf{O que mede} & \textbf{Gold-set} \\
\midrule
Orquestrador puro & Só \emph{tool calling} do LLM (forecaster fixo
              \code{chronos2}) & 25 perguntas \\
Retrospectivo puro & Walk-forward backtest dos forecasters sobre
              séries SINAM isoladas & 30 séries \\
\textbf{Integrado Retrospectivo} & \textbf{Cadeia LLM $\times$ forecaster
              sobre prompts de análise do passado} & \textbf{20 perguntas} \\
\bottomrule
\end{tabularx}
\caption{As três modalidades. As duas primeiras são \emph{diagnósticas}
(isolam componentes). A terceira é \emph{decisional} --- é a que
recomendamos para escolher a combinação que vai pra produção.}
\end{table}

\subsection{Por que rodar o Integrado Retrospectivo}

\begin{itemize}[leftmargin=*]
\item O \textbf{Orquestrador puro} pode dar 100\% em tool calling sem
que o forecaster funcione: ele só verifica \emph{qual} tool foi chamada,
não o que ela retorna.
\item O \textbf{Retrospectivo puro} pode mostrar WAPE excelente em série
isolada, mas se o LLM roteia mal a pergunta o usuário nunca chega a ver
esse WAPE.
\item O \textbf{Integrado Retrospectivo} responde com dados a pergunta
de produto: \emph{``a combinação \code{qwen3:8b $\times$ chattime}
aguenta os prompts retrospectivos típicos dentro de 3\,min, ou
\code{qwen3:8b $\times$ chronos2} é mais previsível?''}.
\end{itemize}

% =============================================================================
\section{Gold-set Integrado Retrospectivo}

O conjunto vive em \file{tests/benchmark/perguntas\_integrado.jsonl} ---
20 perguntas \textbf{retrospectivas} (análise de tendência passada),
divididas em \textbf{dois blocos} de 10:
\begin{itemize}[leftmargin=*]
\item \textbf{Bloco \code{i*}} --- registro analítico genérico
(``Faça a análise retrospectiva de...''). Estilo conciso.
\item \textbf{Bloco \code{j*}} --- registro jurídico-forense
(``Promotor de Justiça do MP-SP requer análise...''). Inclui
referências legais (Lei Maria da Penha, Estatuto do Idoso,
Lei 13.344/2016 de tráfico de pessoas, etc.) e papéis profissionais
(MP, Defensoria, Procuradoria, vara especializada, advocacia).
Testa robustez do LLM a \emph{prompts} formais longos.
\end{itemize}
Cada linha é um JSON com a estrutura:

\begin{lstlisting}[language=Python,caption={Schema de um caso do gold-set retrospectivo.}]
{
  "id": "i01",
  "category": "retrospectivo_uf",
  "question": "Faça a análise retrospectiva de violência sexual em Sergipe nos últimos anos",
  "expected_tools":       ["relatorio_uf"],
  "expected_args":        {"relatorio_uf": {"uf": "SE"}},
  "expected_filter_keys": ["violencia_sexual"],
  "needs_wape":           true,
  "max_seconds":          180
}
\end{lstlisting}

\paragraph{Categorias.} \code{retrospectivo\_uf},
\code{retrospectivo\_municipio}, \code{retrospectivo\_brasil},
\code{retrospectivo\_combo}. Permitem decompor o desempenho por
\emph{tipo} de pergunta retrospectiva (tabela \ref{tab:by-cat-result}).

\paragraph{Por que o nome \emph{forecast} foi descartado.} Apesar de o
pipeline interno usar forecasters em \emph{backtest} para diagnosticar
sazonalidade e tendência, o que o usuário pede é sempre uma narrativa
sobre o passado --- nunca uma projeção futura. Manter o prefixo
\code{forecast\_} no \emph{schema} era confuso para quem lê os
resultados, então padronizamos como \code{retrospectivo\_*}.

% =============================================================================
\section{Métrica de avaliação (sete \emph{checks} booleanos)}

Cada execução é avaliada por \code{evaluate\_case()} em
\file{tests/benchmark/runners/run\_orquestrador.py}, que devolve sete
\emph{flags} booleanas:

\begin{table}[h!]
\centering
\small
\begin{tabularx}{\textwidth}{@{}l X@{}}
\toprule
\textbf{Flag} & \textbf{Verifica} \\
\midrule
\code{correct\_chain}      & Sequência de tools chamadas = \code{expected\_tools} \\
\code{args\_ok}            & UF/município extraído bate com \code{expected\_args} \\
\code{filter\_keys\_ok}    & Filtro de tipo de violência detectado corretamente \\
\code{gold\_kw\_present}   & \emph{Keywords} âncora aparecem na resposta gerada \\
\code{no\_forbidden\_tools}& Não chamou \emph{tool} de outra modalidade \\
\code{no\_forbidden\_kw}   & Não usou palavra-veto (e.g.\ ``diminuiu'' em série crescente) \\
\code{within\_time}        & $\text{elapsed\_s} \le \text{max\_seconds}$ \\
\bottomrule
\end{tabularx}
\caption{Os sete critérios avaliados por execução. Tudo persiste no
\code{per\_combo.csv} de cada \emph{run}.}
\label{tab:flags}
\end{table}

\subsection{Métricas derivadas}

Computadas em \file{webapp/benchmark/views.py} (função
\code{\_integ\_stats}) e em \file{webapp/benchmark/worker.py}
(\code{\_summarise\_integrado}):

\begin{itemize}[leftmargin=*]
\item \textbf{Acurácia} --- média aritmética das sete \emph{flags}, por
linha:
\begin{equation}
\text{acc}_i = \frac{1}{7} \sum_{k=1}^{7} \mathbf{1}[\text{flag}_k^{(i)} = \text{True}]
\end{equation}

\item \textbf{Strict correctness} --- conjunção das três \emph{hard
checks}:
\begin{equation}
\text{strict}_i = \text{correct\_chain}_i \wedge \text{args\_ok}_i \wedge \text{filter\_keys\_ok}_i
\end{equation}

\item \textbf{Full pass} --- todas as sete flags \texttt{True}.

\item \textbf{Latências} --- p50 e p95 do \code{elapsed\_s} dentro do
combo.

\item \textbf{Score} --- combinação linear ponderada (seção \ref{sec:score}).
\end{itemize}

% =============================================================================
\section{Score ponderado}
\label{sec:score}

Para cada par $\langle \text{orq}, \text{fc} \rangle$, define-se:

\begin{equation}
\boxed{
\text{score} \;=\; 0{,}50 \cdot \text{chain\_ok}_\% \;+\; 0{,}20 \cdot \text{args\_ok}_\%
              \;+\; 0{,}15 \cdot \text{filter\_keys}_\% \;+\; 0{,}15 \cdot \text{within\_time}_\%
}
\label{eq:score}
\end{equation}

\paragraph{Justificativa dos pesos.}

\begin{itemize}[leftmargin=*]
\item \textbf{Chain (50\%)} é o erro mais catastrófico --- chamar a
\emph{tool} errada destrói a resposta inteira, independente do que vem
depois.
\item \textbf{Args (20\%)} é parametrização: pega a UF certa mas mira no
filtro errado. Ruim, mas recuperável.
\item \textbf{Filter keys (15\%)} é o tipo de violência detectado;
parametrização secundária.
\item \textbf{Within time (15\%)} pune combos que tecnicamente
\emph{funcionam} mas seriam inaceitáveis em \emph{chat} interativo. É o
peso que reprova o ChatTime apesar da acurácia razoável.
\end{itemize}

% =============================================================================
\section{Produto cartesiano e execução}

O \emph{runner} em \file{webapp/benchmark/worker.py} (função
\code{\_run\_integrado}) constrói o produto cartesiano:

\begin{equation}
\mathcal{E} \;=\; \mathcal{O} \times \mathcal{F} \times \mathcal{C}
\end{equation}

onde $\mathcal{O}$ é o conjunto de orquestradores selecionados,
$\mathcal{F}$ o de forecasters e $\mathcal{C}$ o gold-set
(\(|\mathcal{C}| = 20\); 10 \emph{retrospectivos} + 10 \emph{jurídicos}).
Para cada $(o, f, c) \in \mathcal{E}$ é
disparada uma chamada da cadeia inteira; o resultado vai pra
\file{tests/benchmark/results/$\langle$date$\rangle$\_integ\_run$\langle$id$\rangle$/per\_combo.csv}.

\paragraph{Latência total esperada.} Estimativa:

\begin{equation}
T \;\approx\; |\mathcal{O}| \cdot |\mathcal{F}| \cdot |\mathcal{C}| \cdot \bar{t}_\text{call}
\end{equation}

com $\bar{t}_\text{call} \approx 30\,$s para Qwen3 + Chronos--2 e
$\approx 75\,$s incluindo ChatTime.

% =============================================================================
\section{Implementação --- mapa de arquivos}

\begin{table}[h!]
\centering
\small
\renewcommand{\arraystretch}{1.15}
\begin{tabularx}{\textwidth}{@{}l X@{}}
\toprule
\textbf{Arquivo} & \textbf{Papel} \\
\midrule
\file{tests/benchmark/perguntas\_integrado.jsonl} & 20 perguntas-gold (10 retrospectivas \code{i*} + 10 jurídicas \code{j*}) com schema do §4 \\
\file{tests/benchmark/runners/run\_orquestrador.py} & \code{evaluate\_case()} aplica os 7 \emph{checks} \\
\file{webapp/benchmark/models.py} & \code{BenchmarkRun} (kind, status, progresso, summary) \\
\file{webapp/benchmark/worker.py} & \code{\_run\_integrado()} produto cartesiano; \code{\_summarise\_integrado()} agrega \\
\file{webapp/benchmark/views.py} & \code{run\_new\_integrado}, \code{run\_quick\_integrado}, \code{run\_detail}, \code{\_integ\_stats} \\
\file{webapp/templates/benchmark/run\_new\_integrado.html} & UI: seleção de orq × fc, prévia do cartesiano, estimativa \\
\file{webapp/templates/benchmark/run\_status.html} & UI: \emph{progress bar} + \emph{log} ao vivo (polling a cada 2\,s) \\
\file{webapp/templates/benchmark/run\_detail.html} & UI: KPIs, 4 gráficos Chart.js, \emph{ranking} dos pares \\
\file{webapp/chat/orchestrator.py} & Loop de \emph{tool calling} (OpenAI-compatible) \\
\file{webapp/chat/tools.py} & \code{relatorio\_uf}, \code{relatorio\_municipio}, \code{relatorio\_brasil} \\
\file{src/series_models/\{chronos2,chattime,baselines,random\_forest,remote\_api\}\_adapter.py} &
  Adaptadores de \emph{forecaster} chamados pelas \emph{tools} \\
\bottomrule
\end{tabularx}
\caption{Onde cada peça mora.}
\end{table}

% =============================================================================
\section{Resultado real --- \texttt{2026-06-05\_integ\_run3}}

Combo testado: \emph{qwen3:8b $\times$ chattime}, 10 perguntas (essa
execução é anterior à inclusão do bloco jurídico de junho/26; rodadas
futuras incluem os 20 itens).

\begin{table}[h!]
\centering
\begin{tabular}{@{}l r@{}}
\toprule
\textbf{Métrica} & \textbf{Valor} \\
\midrule
Acurácia (média dos 7 \emph{checks}) & 74{,}3\,\% \\
\code{correct\_chain} \%               & 100{,}0\,\% \\
\code{args\_ok} \%                     & 90{,}0\,\% \\
\code{filter\_keys\_ok} \%             & 70{,}0\,\% \\
\code{within\_time} \%                 & 0{,}0\,\% \\
Strict correctness                    & 30{,}0\,\% \\
Full pass                             & 0{,}0\,\% \\
p50 latência                          & 303{,}2\,s \\
p95 latência                          & 328{,}8\,s \\
\textbf{Score} (eq.~\ref{eq:score})    & \textbf{69{,}0} \\
\bottomrule
\end{tabular}
\caption{Métricas agregadas do \emph{run} \code{integ\_run3}.}
\end{table}

\paragraph{Leitura.} O LLM acertou \emph{tool} em 100\% dos casos e
parâmetros em 90\%. O ponto fraco é o tempo: \code{within\_time} = 0\%
significa que \emph{nenhuma} execução cumpriu o teto de 180\,s --- todas
as chamadas com ChatTime estouraram. Em produção, isso reprova o
ChatTime independentemente da qualidade da previsão.

\begin{table}[h!]
\centering
\small
\begin{tabular}{@{}l r r r r@{}}
\toprule
\textbf{Categoria} & \textbf{N} & \textbf{Chain OK} & \textbf{Args OK} & \textbf{Within time} \\
\midrule
\code{retrospectivo\_uf}        & 5 & 100\% &  80\% & 0\% \\
\code{retrospectivo\_brasil}    & 2 & 100\% & 100\% & 0\% \\
\code{retrospectivo\_municipio} & 2 & 100\% & 100\% & 0\% \\
\code{retrospectivo\_combo}     & 1 & 100\% & 100\% & 0\% \\
\bottomrule
\end{tabular}
\caption{Desempenho por categoria de pergunta retrospectiva
(\emph{breakdown} calculado em \code{\_integ\_stats}).}
\label{tab:by-cat-result}
\end{table}

% =============================================================================
\section{Visualizações}

A página \file{run\_detail.html} renderiza quatro gráficos via Chart.js
em \emph{grid} 2$\times$2:

\begin{enumerate}[leftmargin=*]
\item \textbf{Score por combinação} --- barras coloridas, escala 0--100.
\item \textbf{Perfil de qualidade (radar)} --- cada combo num polígono
nos sete critérios da tabela \ref{tab:flags}. Quanto mais externo,
melhor.
\item \textbf{Latência (p50 vs.\ p95)} --- barras agrupadas; revela
combos com \emph{long tail}.
\item \textbf{Por categoria} --- barras agrupadas mostrando em que tipo
de pergunta o pipeline falha.
\end{enumerate}

% =============================================================================
\section{Reprodução}

\subsection{One-click}

Pelo navegador, em \url{http://127.0.0.1:8000/benchmark/}, basta
clicar no botão \emph{\texttt{Iniciar benchmark automático}}. Ele cria
um \code{BenchmarkRun} com defaults \emph{qwen3:8b $\times$ chronos2} e
redireciona para a tela de status com \emph{progress bar} ao vivo.

\subsection{Customizado}

Em \url{http://127.0.0.1:8000/benchmark/nova/integrado/}:

\begin{enumerate}[leftmargin=*]
\item Selecionar 1 ou mais orquestradores (Ollama local ou APIs
externas cadastradas).
\item Selecionar 1 ou mais \emph{forecasters} (locais ou remotos).
\item Definir \emph{label}, \emph{temperature}.
\item Submeter --- a execução roda em \emph{background} numa thread
\emph{daemon}.
\end{enumerate}

\subsection{Linha de comando}

\begin{lstlisting}[language=bash]
# Linux / macOS
source venv/bin/activate
python tests/benchmark/runners/run_orquestrador.py
\end{lstlisting}

\begin{lstlisting}[language=bash]
# Windows (PowerShell)
.\venv\Scripts\Activate.ps1
python tests\benchmark\runners\run_orquestrador.py
\end{lstlisting}

% =============================================================================
\section{Limitações e trabalho futuro}

\begin{itemize}[leftmargin=*]
\item \textbf{Eval humano ausente.} O protocolo prescritivo
(\file{docs/benchmark\_protocol.md}, §3.1) recomenda Bradley--Terry com
dois avaliadores. Nesta execução, só métricas automáticas foram
computadas.

\item \textbf{Sem \emph{confidence interval}.} O score é pontual; um
\emph{bootstrap} de 1000 amostras daria IC 95\% e permitiria declarar
vencedores com significância estatística.

\item \textbf{Gold-set ainda pequeno.} 20 perguntas (10 retrospectivas + 10 jurídicas) são suficientes para
\emph{sanity check} mas insuficientes para conclusão de produção.
Recomenda-se expandir para $\ge$ 50 cobrindo mais combinações de UF $\times$
tipo de violência $\times$ horizonte temporal.

\item \textbf{Acoplamento das categorias.} A tabela
\ref{tab:by-cat-result} está desequilibrada (5 \code{forecast\_uf},
2 outras, 1 \code{combo}). Balancear o gold-set por categoria.

\item \textbf{Gatilho de adoção formal.} O documento prescritivo (§7)
define quando \emph{trocar} o modelo em produção; o pipeline atual
calcula o score mas não aplica o gatilho automaticamente.
\end{itemize}

% =============================================================================
\section{Conclusão}

O Benchmark Integrado Retrospectivo dá uma resposta \emph{conjunta}
sobre a qualidade da cadeia para o caso de uso real do produto ---
análise de tendência do passado --- com sete checagens objetivas e um
\emph{score} ponderado interpretável. A execução em
\code{2026-06-05\_integ\_run3} mostrou que o gargalo do par
\code{qwen3:8b $\times$ chattime} \emph{não é qualidade} (acurácia 74\%,
\code{chain\_ok} 100\%) mas \emph{latência}: p95 = 329\,s, 100\% dos
casos extrapolam o teto de 180\,s. Em \emph{chat} retrospectivo
interativo, esse tempo é proibitivo. A recomendação prática é manter
\code{chronos2} como \emph{forecaster} padrão (resposta sub-segundo,
qualidade equivalente em retrospectivo) e usar ChatTime apenas em
contextos \emph{batch} --- relatórios noturnos, exportações --- onde o
tempo de execução não pesa.

\bigskip\hrule\smallskip
\noindent\textit{Documento gerado a partir do código em
\code{commit}\textsuperscript{HEAD} do branch \code{main}.
Para a versão prescritiva completa, ver
\file{docs/benchmark\_protocol.md}.}

\end{document}
